\documentclass[11pt,a4paper]{article}

\usepackage[margin=1in]{geometry}
\usepackage{graphicx}
\usepackage{booktabs}
\usepackage{caption}
\usepackage{subcaption}
\usepackage{amsmath, amssymb}
\usepackage{xcolor}
\usepackage[hidelinks, colorlinks=true, linkcolor=blue!50!black,
            citecolor=blue!50!black, urlcolor=blue!60!black]{hyperref}
\usepackage[round, authoryear]{natbib}
\usepackage{titlesec}
\usepackage[protrusion=true,expansion=false]{microtype}
\usepackage{enumitem}
\usepackage{float}

\graphicspath{{../figures/}{../figures/phase0/}{../figures/phase1/}%
              {../figures/phase2/}{../figures/phase3/}{../figures/phase4/}%
              {../figures/phase5/}}

\titleformat{\section}{\large\bfseries}{\thesection.}{0.6em}{}
\titleformat{\subsection}{\normalsize\bfseries}{\thesubsection}{0.6em}{}
\setlength{\parskip}{0.4em}
\setlength{\parindent}{0pt}

\begin{document}

% ---- Custom title block with IITR logo ----
\begin{center}
  \begin{minipage}{0.13\textwidth}
    
\includegraphics[width=\linewidth]{IITR_logo.png}
  \end{minipage}\hfill
  \begin{minipage}{0.83\textwidth}
    \raggedright
    {\small Indian Institute of Technology Roorkee \\
     \textbf{DAL-311 --- Network Science}}\\[6pt]
    {\Large \textbf{Trade-Network Resilience and the Strait of Hormuz:}}\\
    {\large \textbf{A Multi-Layer Exposure Analysis of Five Commodity
                    Networks (2020--2024)}}
  \end{minipage}
\end{center}

\bigskip
\begin{center}
  \textbf{Group members}\\[4pt]
  \begin{tabular}{ll}
    Preetham Gunta      & Vineel Krishna Manchala \\
    Laluprasad Ramavath & Prabhath Kolli          \\
  \end{tabular}
  \\[6pt]
  \textit{May 2026}
\end{center}
\rule{\textwidth}{0.4pt}\medskip

\begin{abstract}
The Strait of Hormuz carries roughly a fifth of the world's
seaborne energy trade, and a closure or partial blockade would propagate
asymmetrically across commodity-specific trade networks. We assemble
six years (2020--2025) of bilateral monthly trade flows, tanker-capacity
indicators and benchmark prices to build five directed weighted commodity
networks --- wheat, ammonia, urea, LPG propane, and LPG butane --- and
measure their exposure to a Hormuz disruption. Using community
detection, weighted topological metrics and targeted-attack simulations,
we show that LPG butane is the most fragile layer: removing Hormuz
exporters wipes out 48.5\,\% of trade volume and 32.4\,\% of weighted
global efficiency in 2024. A bivariate Vector AutoRegression and a
Spatial Durbin panel model identify a small but consistent network-
contagion coefficient ($\rho \approx 0.11$) and a statistically
significant Granger link from Hormuz tanker capacity to ammonia prices.
A scenario engine combining structural disruption with the
spatial-econometric multipliers projects that under a six-month
realignment scenario, India alone faces an annualised import-value loss
of \$7.6 billion --- driven entirely by its $\approx\!97\,\%$ Hormuz
dependence on LPG. We discuss the methodological limits of the
econometric channel and the policy implications of LPG concentration.
\end{abstract}

% =====================================================================
\section{Introduction and problem statement}
% =====================================================================

The Strait of Hormuz is the maritime corridor between the Persian Gulf
and the Gulf of Oman, through which an estimated 17--20\,\% of global
oil-equivalent energy and over a quarter of global LPG trade transits
\citep{eia_hormuz}. Repeated geopolitical episodes (the 2019 tanker
seizures, the 2024 Iran--Israel exchange, recurring threats to ``close''
the strait) have repeatedly raised the question: \emph{how exposed is
each commodity-specific trade network to a Hormuz disruption, and how
would a shock propagate across countries?}

\paragraph{Research questions.} This project tackles five concrete
questions:
\begin{enumerate}[leftmargin=2em, itemsep=0.1em]
  \item \textbf{Structure.} Are the five commodity networks
        topologically similar, or do they have distinct exposure
        signatures to Hormuz exporters?
  \item \textbf{Exposure.} Which countries are most exposed, and is
        ``exposure'' best captured by direct bilateral share, by
        global centrality, or by an information-flow measure?
  \item \textbf{Resilience.} How much trade volume and connectivity is
        lost when Hormuz exporter rows are removed from each network?
  \item \textbf{Dynamics.} Does Hormuz tanker capacity Granger-cause
        commodity price movements, and does shock propagation across
        the network amplify imports' co-movement?
  \item \textbf{Scenarios.} Under a 14-day shock, a 30--45 day
        rupture, and a 6-month realignment of trade routes, what is
        the projected country-level dollar impact on imports?
\end{enumerate}

\paragraph{Why a multi-layer treatment.} Treating ``trade'' as a single
flow understates the very different role Hormuz plays for, e.g.,
LPG (where Gulf states are dominant exporters) versus wheat (where they
are net importers). The five layers chosen here separate energy
(LPG~butane, LPG~propane), agriculture (wheat) and fertilizer
(ammonia, urea) so that the structural effects of a Hormuz disruption
are not averaged out across heterogeneous flows.

% =====================================================================
\section{Data sources and collection}
% =====================================================================

The pipeline draws on five public, machine-readable databases:

\begin{itemize}[leftmargin=1.4em, itemsep=0.15em]
  \item \textbf{UN Comtrade Plus} \citep{comtrade}: bilateral monthly
        trade flows by HS code. We extract HS~1001 (wheat),
        HS~2814 (ammonia), HS~3102\,10 (urea), HS~2711\,12 (LPG
        propane) and HS~2711\,13 (LPG butane) for 2020--2025
        (204\,019 reporter-partner-month observations).
  \item \textbf{IMF PortWatch} \citep{portwatch}: daily counts and
        deadweight tonnage (DWT) of vessels transiting global
        chokepoints. We use the Hormuz daily tanker series, aggregated
        to monthly $n_\text{tanker}$ and capacity $\text{DWT}$.
  \item \textbf{World Bank Pink Sheet} \citep{worldbank_cmo}: global
        benchmark prices --- US~HRW wheat, urea, DAP (used as ammonia
        proxy), and LNG~Japan (used as butane proxy). Monthly,
        2020-01 to 2026-02.
  \item \textbf{FRED} \citep{fred}: daily Mont~Belvieu propane spot,
        averaged to monthly. Pre-2021-03 values are backfilled from a
        scaled LNG~Japan index (overlap-mean calibration on
        2021-03..2021-12).
  \item \textbf{JODI World Database} \citep{jodi}: country-level NGL
        production, imports, exports and stocks (used for
        reserve-threshold calibration of Phase~1 robustness checks).
\end{itemize}

% =====================================================================
\section{Data preprocessing}
% =====================================================================

Each raw source is normalised into a tidy CSV under
\texttt{Processed\_Data/}. The main interventions are:

\begin{itemize}[leftmargin=1.4em, itemsep=0.15em]
  \item \emph{Comtrade.} A 47-column / 46-value header shift in the
        upstream CSV is corrected by dropping a leading
        \texttt{typeCode} field and a trailing pad. We remove
        964 self-loops (a country reporting trade with itself) and
        21 constant or unused columns. World-totals and bilateral
        rows are split into two files; HS codes are mapped to short
        layer labels.
  \item \emph{Pink Sheet.} The two-row Excel header band is flattened
        (commodity name + unit on the prices sheet, nested category
        on the indices sheet); the multilingual missing-value markers
        \texttt{\dots} and the cp1252 replacement glyph are coerced
        to \texttt{NaN}.
  \item \emph{FRED propane.} A 2020-01..2021-02 coverage gap is
        backfilled from the LNG-Japan series scaled to its 2021
        overlap mean, so that log-differences remain continuous.
  \item \emph{PortWatch.} Daily UTC timestamps are parsed; zero-traffic
        days (legitimate, not missing) are preserved; daily counts
        and DWT are summed to monthly with a \texttt{n\_days}
        quality flag.
  \item \emph{JODI.} Seven yearly files are concatenated, filtered to
        the \texttt{KTONS} unit to remove a 5-unit redundancy, and
        pivoted to a wide NGL panel covering 65--72 reporters with
        89.6--96.6\,\% cell completeness on the four flows used for
        reserve calibration.
\end{itemize}

The cleaned monthly Comtrade panel covers 153 reporter countries and
five layers across 72 months (2020-01..2025-12). Country ISO3 codes are
used throughout; non-ISO reporters (\texttt{S19} ``Other Asia'') are
dropped.

% =====================================================================
\section{Methodology}
% =====================================================================

The pipeline is organised in five phases. Methods that are well-known
or borrowed from the literature are stated briefly with a citation; the
specific metrics adapted to the trade-network context are explained in
more detail.

\subsection{Phase 0 --- Trade matrix construction}

For each (year, layer) we build a $N \!\times\! N$ matrix
$W_{y,L}$ with $W_{ij}$ equal to the USD value exported from country
$i$ to country $j$, where $N = 153$. We hold the country index fixed so
that all matrices share a common axis. The binary support
$A_{y,L} = \mathbf{1}\{W_{y,L} > 0\}$ is also stored. Bilateral flows
are reconciled with World totals using a single-pass
RAS-style row/column rebalance to keep marginal sums consistent under
small reporter / partner mismatches; the 2020 \texttt{DEU} HS\,1001
sanity check matches to floating-point tolerance.

\subsection{Phase 1 --- Composite exposure score}

For each (country, year, layer) we compute eight exposure metrics and
combine them into a single $z$-score composite (rescaled to
[0, 1]). The metrics are:

\begin{itemize}[leftmargin=1.4em, itemsep=0.15em]
  \item \textbf{DE} (\emph{Direct Exposure}) --- the share of country
        $i$'s imports in layer $L$ that come from Hormuz exporters.
        A simple bilateral-trade ratio.
  \item \textbf{SR} (\emph{Source Redundancy}) --- the
        Herfindahl-Hirschman concentration of country $i$'s import
        sources; high SR means imports come from few partners and so
        the country has little fallback capacity.
  \item \textbf{Betweenness centrality} --- the fraction of weighted
        shortest paths in the network that pass through node $i$.
        High betweenness means $i$ sits on many trade routes.
  \item \textbf{Katz centrality} --- a recursive count of all walks
        ending at $i$, attenuated by a factor $\alpha$. It rewards
        nodes that are reached by many paths, including indirect ones.
  \item \textbf{PPR} (\emph{Personalised PageRank}) --- the steady-state
        probability that a random walker, restarted at the Hormuz set,
        visits node $i$. A natural ``influence-from-Hormuz'' score.
  \item \textbf{Leontief multiplier} --- the diagonal of
        $(I-\beta W^\top)^{-1}$, capturing a country's exposure to
        upstream supply through all rounds of intermediate trade
        (the input-output amplification).
  \item \textbf{DebtRank} \citep{battiston2012} --- a recursive
        ``shock from Hormuz'' diffusion that respects each node's
        absorption capacity, originally developed for systemic
        financial risk.
  \item \textbf{PIVI} (\emph{Path-Integrated Vulnerability Index}) ---
        a compound metric that integrates exposure across all
        directed paths from Hormuz to the focal country, weighted by
        path-attenuation.
\end{itemize}

The eight metrics are standardised ($z$-scored within each year and
layer), averaged, and rescaled to the $[0, 1]$ \texttt{composite\_01}
score (Phase~1.2). The eight metrics are positively but not perfectly
correlated (Figure~\ref{fig:metric-corr}), so the composite is more
than a relabelling of any single index.

\subsection{Phase 2 --- Network topology and resilience}

\textbf{Community detection.} We run InfoMap
\citep{rosvall2008} --- an information-theoretic community-detection
algorithm that minimises the description length of a random walk on
the network. Communities are reported per (year, layer); the
\emph{primary Hormuz community} is defined as the single module
containing the most Hormuz members (ties broken by smaller size). This
fixes a labelling-instability issue under InfoMap's stochastic seed
and gives a stable measure of bloc cohesion.

\textbf{Topology metrics.} For each network we compute density,
overall reciprocity \citep{garlaschelli2004},
average undirected clustering, the size of the largest weakly /
strongly connected components (LWCC, LSCC) and weighted
\emph{global efficiency} \citep{latora2001} on cost~$=1/W$.

\textbf{Targeted attack.} The Hormuz attack zeroes out all
out-going edges from the seven Hormuz countries
$\{$ARE, BHR, IRN, IRQ, KWT, QAT, SAU$\}$ and reports
flow loss, importers cut to zero, the relative LWCC retained, and the
relative drop in global efficiency. This is the standard targeted
node-removal scheme of \citet{albert2000} adapted to weighted directed
graphs.

\subsection{Phase 3 --- Panel econometrics}

\textbf{Bivariate VAR.} For each layer we fit a Vector AutoRegression
\citep{sims1980} on
$[\,\Delta\!\log\text{price}_t,\, \text{DWT}_t\,]$, with lag
order chosen by BIC (max 4) and Cholesky identification ordering DWT
first. Granger causality \citep{granger1969} is tested in the
``DWT $\rightarrow$ price'' direction, and a 12-month orthogonalised
impulse-response function (IRF) is reported with asymptotic 95\,\%
confidence bands.

\textbf{Panel regression.} For each layer we estimate
\[
  \log \text{import}_{it}
   \;=\; \alpha_i \;+\; \alpha_t
   \;+\; \beta_1 \log p_t
   \;+\; \beta_2 \,\text{DWT}_t
   \;+\; \gamma_3 \,\text{DE}_i \!\times\! \text{DWT}_t
   \;+\; \varepsilon_{it}
\]
with country and time fixed effects. The interaction
$\gamma_3$ is the parameter of substantive interest --- it tests
whether more-exposed countries (higher DE) co-move with Hormuz
capacity differently from less-exposed ones.

\textbf{Spatial Durbin Model.} The annual analogue uses
\texttt{spreg.Panel\_FE\_Lag} with W~=~row-normalised 2020 trade-share
matrix \citep{lesage2009, rey2010}:
\[
  \log\text{import}_{it} = \rho\,W \log\text{import}_{t}
   + X_{it}\boldsymbol\beta + \mu_i + \varepsilon_{it}.
\]
The spatial autoregressive coefficient $\rho$ captures
\emph{network contagion}: how much a country's imports co-move with
those of its weighted trade partners after controlling for direct
exposure and global price. We complement this with an SLX (Spatial
Lag of $X$) specification that estimates the indirect / spillover
effect $\theta$ of price and DE without the simultaneity of the lag
in $Y$.

\subsection{Phase 4 --- Scenario simulation engine}

Three Hormuz disruption scenarios are simulated:

\smallskip
\begin{tabular}{ll}
\textbf{S1 --- 14-day transient shock} &
   $\Delta\!\text{DWT}_t = -0.5$ for $t = 1$, then 0. \\[2pt]
\textbf{S2 --- 30--45 day rupture} &
   $\Delta\!\text{DWT}_t = (-0.9, -0.9, -0.4, 0, \dots)$. \\[2pt]
\textbf{S3 --- 6-month persistent realignment} &
   $\Delta\!\text{DWT}_t = -0.7$ for $t = 1, \dots, 6$, tapered after.
\end{tabular}
\smallskip

For each scenario, country, layer and month we sum two channels:

\begin{enumerate}[leftmargin=1.4em, itemsep=0.1em]
  \item \emph{Structural channel.} $\log(1 - \text{DE}_i \cdot
        \max(-\Delta\!\text{DWT}_t, 0))$. This captures the trade-
        network supply disruption in proportion to the country's
        Hormuz dependence.
  \item \emph{Price channel.} The long-run cumulative IRF of
        $\Delta\!\log p$ to a Hormuz DWT shock, scaled by the SDM
        $\beta_{\log p}$ coefficient, gives a per-period log-import
        deviation at a given DWT level.
\end{enumerate}

The two channels are added and multiplied by the SDM Leontief
amplifier $1 / (1 - \rho)$. To avoid the explosive accumulation
that arises from convolving the IRF over a sustained shock, the
price channel uses the long-run IRF as a per-period steady-state
elasticity. Cumulative monthly log-changes are normalised to an
average-monthly metric (divide by 12) before exponentiating to a
percentage. Annual dollar impact is the average percentage change
applied to 2024 baseline imports.

% =====================================================================
\section{Results}
% =====================================================================

\subsection{The five commodity networks at a glance}

Figure~\ref{fig:network-panel} shows the top-60 weighted edges of each
2024 commodity network. Across all panels, \emph{nodes} are countries
(ISO3); \emph{edges} are directed exporter $\rightarrow$ importer
flows; \emph{node size} scales with weighted in+out degree (a country's
total trade involvement); \emph{node colour} flags Hormuz exporters in
red. Edge width and opacity are proportional to trade weight.

\begin{figure}[H]
  \centering
  \includegraphics[width=\textwidth]{fig5_network_panel_2024}
  \caption{Top-60-edge bilateral trade networks for the five
  commodity layers, 2024. Hormuz exporter states (red) are central
  hubs in both LPG layers and dominate ammonia/urea flows; in wheat
  they appear only as importers (small red satellites).}
  \label{fig:network-panel}
\end{figure}

The visual summary mirrors the global statistics: in LPG butane and
LPG propane the red Hormuz nodes are large and well connected; in
wheat they are small importer nodes at the periphery. The ammonia
and urea networks are intermediate, with Gulf exporters playing a
substantial role.

\begin{figure}[H]
  \centering
  \includegraphics[width=0.95\textwidth]{fig5_network_lpg_butane_2024}
  \caption{LPG butane bilateral trade network (2024), top 90 edges.
  Saudi Arabia (SAU), United Arab Emirates (ARE) and Qatar (QAT) are
  the dominant export hubs; India (IND), Japan (JPN) and China
  (CHN) are the largest importer hubs.}
  \label{fig:network-lpg}
\end{figure}

\subsection{Hormuz share of trade and exposure landscape}

\begin{figure}[H]
  \centering
  \includegraphics[width=0.85\textwidth]{fig0_hormuz_share}
  \caption{Hormuz exporter share of global trade volume by layer and
  year. LPG butane is the most concentrated (35--50\,\%), wheat the
  least (under 1\,\%).}
  \label{fig:hormuz-share}
\end{figure}

The Phase~1 composite exposure (Figure~\ref{fig:composite-heatmap})
combines all eight metrics into a single $[0, 1]$ score. The eight
underlying metrics correlate positively but capture distinct
information (Figure~\ref{fig:metric-corr}): DE and PPR are highly
correlated, while betweenness and DebtRank pick up complementary
structural roles.

\begin{figure}[H]
  \centering
  \begin{subfigure}{0.48\textwidth}
    \includegraphics[width=\linewidth]{fig1_composite_heatmap_2024}
    \caption{2024 composite exposure heatmap.}
    \label{fig:composite-heatmap}
  \end{subfigure}\hfill
  \begin{subfigure}{0.48\textwidth}
    \includegraphics[width=\linewidth]{fig1_metric_correlation}
    \caption{Cross-metric correlation across 5 layers and 5 years.}
    \label{fig:metric-corr}
  \end{subfigure}
  \caption{Phase~1 composite exposure (left) is concentrated on
  South and East Asian importers in the LPG layers. The eight
  underlying metrics (right) correlate positively but are not
  redundant: bilateral DE and the random-walk PPR cluster together,
  while betweenness and DebtRank pick up complementary structural
  roles.}
\end{figure}

The composite ranking confirms what visual inspection of the network
suggests, but quantifies it: India tops the LPG butane and LPG
propane ranks with composite scores of 1.00, driven by 96--98\,\%
direct dependence (DE) on Hormuz exporters and the largest absolute
import volumes in the world. Table~\ref{tab:top-composite} lists the
top-7 most-exposed countries in each LPG layer in 2024.

\begin{table}[H]
\centering
\small
\caption{Top-7 most-exposed countries in 2024 by composite score.}
\label{tab:top-composite}
\begin{tabular}{lcc@{\hskip 1.2cm}lcc}
\toprule
\multicolumn{3}{c}{\textbf{LPG butane}} & \multicolumn{3}{c}{\textbf{LPG propane}} \\
\cmidrule(r){1-3}\cmidrule(l){4-6}
Country & DE & composite & Country & DE & composite \\
\midrule
India       & 0.98 & 1.00 & India       & 0.97 & 1.00 \\
China       & 0.61 & 0.77 & China       & 0.28 & 0.93 \\
Hong Kong   & 0.00 & 0.61 & Myanmar     & 0.00 & 0.69 \\
Thailand    & 0.81 & 0.51 & Thailand    & 0.91 & 0.52 \\
Malaysia    & 0.59 & 0.50 & Philippines & 0.90 & 0.50 \\
Philippines & 0.81 & 0.49 & Yemen       & 1.00 & 0.48 \\
Singapore   & 0.81 & 0.46 & Oman        & 1.00 & 0.48 \\
\bottomrule
\end{tabular}
\end{table}

The presence of Hong Kong and Myanmar with DE = 0 but high composite
score is a useful sanity check on the methodology: these countries
do not buy LPG \emph{directly} from Hormuz states but sit on the
trans-shipment routes that re-distribute Gulf cargoes through
Singapore and the South-China Sea --- a structural exposure that the
betweenness, Katz and PPR components of the composite correctly pick
up even when the direct bilateral share is zero.

\begin{figure}[H]
  \centering
  \begin{subfigure}{0.49\textwidth}
    \includegraphics[width=\linewidth]{fig1_top_countries_lpg_butane_2024}
    \caption{Top countries by composite exposure (2024).}
  \end{subfigure}\hfill
  \begin{subfigure}{0.49\textwidth}
    \includegraphics[width=\linewidth]{fig1_top5_timeseries_lpg_butane}
    \caption{Composite trajectory 2020--2024.}
  \end{subfigure}
  \caption{LPG butane exposure: India sits well above all peers in
  every year of the panel; China rises noticeably from 2022 onward
  as its butane demand outgrows domestic supply. The top-5 ranks are
  stable across the five years --- the exposure landscape is a
  structural feature of the network, not a year-specific artefact.}
  \label{fig:lpg-top-countries}
\end{figure}

\subsection{Network topology over time}

Topological metrics are stable across 2020--2024
(Table~\ref{tab:topology}, Figure~\ref{fig:topology-grid}). All five
networks have a single weakly-connected component
(LWCC $\approx 1$), low density (5--8\,\%) and reciprocity
0.25--0.34. Wheat has the highest reciprocity (many countries are
both exporter and importer), urea the densest network.

\begin{table}[H]
\centering
\small
\caption{Mean topology metrics across 2020--2024 by layer.}
\label{tab:topology}
\begin{tabular}{lccccc}
\toprule
Layer & Density & Reciprocity & Avg.\ clustering & LWCC frac.\ &
   Global eff.\ ($\times 10^6$) \\
\midrule
Wheat       & 0.065 & 0.342 & 0.484 & 1.000 & 3.17 \\
Ammonia     & 0.062 & 0.289 & 0.554 & 1.000 & 0.81 \\
Urea        & 0.084 & 0.335 & 0.507 & 0.997 & 1.69 \\
LPG propane & 0.048 & 0.254 & 0.424 & 1.000 & 3.65 \\
LPG butane  & 0.056 & 0.273 & 0.457 & 1.000 & 1.73 \\
\bottomrule
\end{tabular}
\end{table}

\begin{figure}[H]
  \centering
  \includegraphics[width=0.92\textwidth]{fig2_topology_grid}
  \caption{Density, reciprocity, average clustering and LWCC fraction
  by year and layer. The five networks evolve smoothly with no
  structural break across the modelling window.}
  \label{fig:topology-grid}
\end{figure}

A few patterns are worth pointing out. \emph{Reciprocity} is highest
in wheat and urea ($\approx\!0.34$): many countries are simultaneously
import\-ers and re-exporters of these commodities, so most edges have
a counter-flow. The two LPG layers have the lowest reciprocity
($0.25$--$0.27$) --- the Gulf-to-Asia LPG trade is overwhelmingly
one-directional, an asymmetry that compounds vulnerability under a
Hormuz disruption. \emph{Average clustering} ranges from 0.42
(LPG propane) to 0.55 (ammonia), with the higher figures indicating
denser local triangles around regional fertilizer hubs (e.g.\ the
Gulf-Russia-Egypt fertilizer triangle).

\paragraph{Community detection.}
InfoMap returns 8--15 communities per layer. The two LPG layers are
the most fragmented (10 and 15 communities respectively), reflecting
multiple regional sub-markets (East Asia, West Africa, Latin
America). The ammonia and urea networks are more centralised
(8 communities), and the wheat layer is split into eight modules
that broadly mirror geographic continents.

\begin{table}[H]
\centering
\small
\caption{2024 community structure and the primary Hormuz bloc.}
\label{tab:communities}
\begin{tabular}{lcccc}
\toprule
Layer & \#\,communities & Modularity Q & Primary bloc size &
   \% of H exports inside \\
\midrule
Wheat       & 8  & 0.50 & 20 & 31\,\% \\
Ammonia     & 8  & 0.57 & 32 & \textbf{88\,\%} \\
Urea        & 8  & 0.41 & 33 & 46\,\% \\
LPG propane & 10 & 0.30 & 13 & 61\,\% \\
LPG butane  & 15 & 0.45 & 19 & \textbf{90\,\%} \\
\bottomrule
\end{tabular}
\end{table}

The most striking column of Table~\ref{tab:communities} is the share
of Hormuz exports that stay inside the primary Hormuz community.
For ammonia and LPG~butane, $\sim\!90\,\%$ of all Hormuz outflow
lands within a single tightly-bound module of $\sim\!20$--$30$
countries --- the cluster of importers that is functionally
\emph{dependent} on the Hormuz route. For wheat, the figure is only
31\,\%, simply because Hormuz states are wheat importers themselves
and their tiny export flows go in many different directions. The
combination of high concentration (ammonia, LPG~butane) with high
Hormuz centrality identifies the layers where a regional shock has
a clean transmission channel through a structurally cohesive bloc.

\begin{figure}[H]
  \centering
  \begin{subfigure}{0.49\textwidth}
    \includegraphics[width=\linewidth]{fig2_community_count}
    \caption{Number of communities (InfoMap, weighted directed)
             per layer over time.}
  \end{subfigure}\hfill
  \begin{subfigure}{0.49\textwidth}
    \includegraphics[width=\linewidth]{fig2_hormuz_bloc_share}
    \caption{Share of Hormuz exports that lands inside the primary
             Hormuz community.}
  \end{subfigure}
  \caption{Community structure is stable year-to-year. The high and
  flat curves for ammonia and LPG butane (right panel) confirm that
  the Gulf-importer bloc is a persistent feature of those networks,
  not a 2024-specific artefact.}
  \label{fig:community-trajectories}
\end{figure}

\subsection{Resilience under Hormuz removal}

The headline resilience finding is the asymmetric impact of removing
Hormuz exporters across layers (Figure~\ref{fig:attack-flow-loss},
Table~\ref{tab:attack}). LPG butane loses 48.5\,\% of trade volume
and 32.4\,\% of weighted global efficiency in 2024; LPG propane and
urea lose 26\,\% and 21\,\%. Wheat is essentially unaffected
(0.13\,\% flow loss), reflecting the Gulf's role as a wheat importer,
not exporter.

\begin{figure}[H]
  \centering
  \begin{subfigure}{0.46\textwidth}
    \includegraphics[width=\linewidth]{fig2_attack_flow_loss}
    \caption{Share of trade volume removed in 2024.}
  \end{subfigure}\hfill
  \begin{subfigure}{0.50\textwidth}
    \includegraphics[width=\linewidth]{fig2_attack_efficiency}
    \caption{Drop in global efficiency by year and layer.}
  \end{subfigure}
  \caption{Targeted Hormuz removal: LPG butane is the most fragile
  layer; wheat the most resilient.}
  \label{fig:attack-flow-loss}
\end{figure}

\begin{table}[H]
\centering
\small
\caption{Hormuz-removal damage (2024) by layer.}
\label{tab:attack}
\begin{tabular}{lcccc}
\toprule
Layer & Flow loss & Eff.\ drop & Importers cut & LWCC retained \\
\midrule
Wheat       & 0.13\,\% & $-3.8$\,\% & 0 & 0.98 \\
Ammonia     & 12.7\,\% & 9.3\,\% & 0 & 0.97 \\
Urea        & 21.3\,\% & 14.2\,\% & 0 & 0.97 \\
LPG propane & 25.6\,\% & 9.3\,\% & 1 & 0.94 \\
LPG butane  & \textbf{48.5\,\%} & \textbf{32.4\,\%} & 0 & 0.96 \\
\bottomrule
\end{tabular}
\end{table}

The wheat \emph{negative} efficiency drop is a network-pruning
artefact: removing Hormuz countries (which are wheat importers,
not exporters) deletes minor re-export edges and shortens average
weighted path lengths --- a benign effect that does not indicate any
welfare gain.

\paragraph{Severity gradient.}
The five layers split cleanly into a three-tier exposure ranking:

\begin{itemize}[leftmargin=1.4em, itemsep=0.1em]
  \item \emph{Critical} --- LPG butane (48.5\,\% volume loss,
        32.4\,\% efficiency drop). The Gulf is the dominant
        butane exporter and the importer ecosystem is tightly
        clustered (88\,\% within-bloc concentration), so the network
        has very little capacity to reroute.
  \item \emph{Severe} --- LPG propane and urea
        (21--26\,\% volume loss, $\approx\!10$--$14$\,\%
        efficiency drop). LPG propane has more diverse exporters
        (US shale gas, Algeria) so the residual network can absorb
        some flow, while urea has high reciprocity allowing
        partial substitution.
  \item \emph{Mild} --- ammonia (12.7\,\% volume loss, 9.3\,\%
        efficiency drop) and wheat (essentially unaffected).
\end{itemize}

\paragraph{Importer-cut count.}
Only one country (in the LPG-propane layer) sees its in-degree drop
to zero under the attack. This is a useful negative result: most
``exposed'' countries do not lose all of their suppliers, they lose
their \emph{cheapest and largest} suppliers. The economic cost is
therefore better captured by the volume-loss and efficiency-drop
metrics than by the binary ``is the country isolated'' indicator.

\subsection{Time-series dynamics: VAR and impulse responses}

Figure~\ref{fig:price-vs-dwt} overlays the global commodity prices
with the normalised Hormuz tanker DWT series. Visually, all five
prices co-move over the 2022--2023 period when DWT is rising;
quantitatively, only \textbf{ammonia} (Granger $p = 0.010$)
shows a statistically significant lead from DWT to price growth at a
1-month horizon (Figure~\ref{fig:irf-ammonia}). The IRF peaks at
$-0.015$ in $\Delta\!\log p$ at $h = 1$, confirming that a positive
shock to Hormuz tanker capacity (more supply) lowers the
fertilizer-proxy price.

\begin{figure}[H]
  \centering
  \includegraphics[width=0.92\textwidth]{fig3_price_vs_DWT}
  \caption{Global commodity prices vs Hormuz tanker DWT, 2020--2024.
  Each panel shows the layer's log-price (left axis, coloured) and
  the normalised Hormuz tanker DWT (right axis, grey dashed).}
  \label{fig:price-vs-dwt}
\end{figure}

\begin{figure}[H]
  \centering
  \includegraphics[width=0.65\textwidth]{fig3_irf_ammonia}
  \caption{Ammonia (DAP-proxy) impulse response: a $-1\sigma$ Hormuz
  DWT shock raises $\Delta\!\log p$ by $\sim 0.015$ at one-month
  horizon. The 95\,\% asymptotic CI does not cross zero at $h = 1$.}
  \label{fig:irf-ammonia}
\end{figure}

\subsection{Spatial network contagion: SDM}

The Spatial Durbin Model identifies a positive spatial-lag
coefficient $\rho \approx 0.11$ for wheat, ammonia and LPG butane,
suggestive of network co-movement of import volumes after
controlling for direct exposure and price
(Figure~\ref{fig:sdm-rho}). The estimates are significant at the
10\,\% level given the small annual panel ($T = 5$), and translate
into a Leontief amplifier $1/(1-\rho) \approx 1.13$ --- a 13\,\%
amplification of any direct shock through the network.

\begin{figure}[H]
  \centering
  \begin{subfigure}{0.46\textwidth}
    \includegraphics[width=\linewidth]{fig3_rho_by_layer}
    \caption{SDM $\rho$ estimates by layer.}
  \end{subfigure}\hfill
  \begin{subfigure}{0.50\textwidth}
    \includegraphics[width=\linewidth]{fig3_sdm_direct_indirect}
    \caption{SLX direct ($\beta$) vs.\ indirect ($\theta$) effects.}
  \end{subfigure}
  \caption{Network contagion is small but consistent in sign for
  the most-traded layers.}
  \label{fig:sdm-rho}
\end{figure}

\subsection{Scenario simulation}

The DWT profiles for the three scenarios and the resulting cross-
layer impacts are shown in Figure~\ref{fig:scenario-comparison}.
The 6-month realignment scenario (S3) projects the largest LPG
disruptions: \$5.7\,billion in lost butane imports and \$4.5\,billion
in lost propane imports globally over the year.

\begin{figure}[H]
  \centering
  \includegraphics[width=\textwidth]{fig4_scenario_comparison}
  \caption{Left: DWT capacity profile for all three scenarios.
  Right: aggregate annualised import-value impact by commodity for
  the two short-horizon scenarios (S1 and S2). The 6-month
  realignment is shown in the heatmap and country-ranking views
  below but is excluded from this aggregate bar chart because the
  IRF-driven price channel is extrapolated beyond its estimation
  horizon under a sustained shock of that length.\
  LPG layers carry essentially the entire negative impact; the small
  positive bars for wheat / ammonia / urea reflect the price channel
  and are an estimation artefact (\S\ref{sec:caveats}).}
  \label{fig:scenario-comparison}
\end{figure}

\paragraph{Why are the wheat / ammonia / urea bars positive?}
The simulation has two channels: a structural channel
($\log(1 - \text{DE}\cdot\delta)$) and a price channel
($\beta_{\log p}\!\cdot\!\Delta\!\log p$). For wheat, ammonia and urea,
\emph{the structural channel is essentially zero for almost every
country in the world} --- only $\sim\!1$--$13\,\%$ of these
commodities transit Hormuz, so the average importer has DE near zero.
The simulated impact is then dominated by the price channel: when
Hormuz capacity falls, prices rise, and the SDM
$\beta_{\log p}\!>\!0$ pushes \emph{nominal} imports up.

This positive co-movement between prices and import volumes is a
real feature of the 2020--2024 estimation window --- it includes the
post-COVID food-security rush (governments stockpiling grain) and
the 2022--23 fertilizer-price spike that triggered emergency
fertilizer subsidies in importing countries. It does not mean these
countries \emph{benefit} from a Hormuz disruption; it means our
econometric channel cannot separate ``defensive stockpiling at high
prices'' from ``ordinary demand response'' on this short panel. We
therefore treat the LPG bars (where the structural channel
dominates) as the policy-relevant signal and the wheat / ammonia /
urea bars as estimation noise.

\begin{figure}[H]
  \centering
  \includegraphics[width=0.78\textwidth]{fig4_heatmap_S3_6month_realignment}
  \caption{S3 (6-month realignment) country~$\times$~layer heatmap of
  average monthly log-import change for the top-40 importers.
  Saturated red cells in the LPG columns identify
  India, Bangladesh, Pakistan and Japan as the worst affected.}
  \label{fig:heatmap-s3}
\end{figure}

The S3 country ranking is dominated by countries with both large
absolute LPG imports and high Hormuz dependence. India is the
single largest absolute loser (Table~\ref{tab:top-losers},
Figure~\ref{fig:top-losers}); its LPG propane and butane DE
exceed 0.96, so a 70\,\% Hormuz capacity reduction sustained for
six months translates into a 46--52\,\% average monthly import
shortfall on those layers.

\begin{table}[H]
\centering
\small
\caption{Top-10 absolute dollar losers under S3 (6-month
realignment). Total dollar = annualised import-value loss across
all five layers.}
\label{tab:top-losers}
\begin{tabular}{lrr}
\toprule
Country & Annual loss (USD M) & \% of 2024 baseline \\
\midrule
India       & $-7\,630$ & $-42$ \\
China       & $-2\,512$ & $-10$ \\
Thailand    & $-572$    & $-17$ \\
Indonesia   & $-340$    & $-5$  \\
Australia   & $-284$    & $-17$ \\
Singapore   & $-85$     & $-26$ \\
Malaysia    & $-79$     & $-8$  \\
Philippines & $-69$     & $-2$  \\
South Africa& $-34$     & $-3$  \\
New Zealand & $-31$     & $-7$  \\
\bottomrule
\end{tabular}
\end{table}

\begin{figure}[H]
  \centering
  \includegraphics[width=0.85\textwidth]{fig4_top_losers}
  \caption{Top-30 most-affected countries under S3, ranked by
  absolute annual dollar impact. India is an outlier driven by
  $\approx\!97\,\%$ Hormuz dependence on LPG.}
  \label{fig:top-losers}
\end{figure}

% =====================================================================
\section{Discussion and limitations}\label{sec:caveats}
% =====================================================================

\paragraph{Robust findings.}
The structural picture is clear and consistent across methods:

\begin{itemize}[leftmargin=1.4em, itemsep=0.1em]
  \item Hormuz removal is asymmetric across commodity layers, with
        LPG butane losing nearly half of global trade volume.
  \item Country-level exposure to Hormuz disruption is concentrated
        in a handful of large LPG-importing economies: India, China,
        Thailand, Indonesia, Australia, Singapore.
  \item The composite exposure score correlates with --- but is not
        identical to --- the simpler bilateral dependence index DE,
        confirming that network position adds information beyond
        the trade ratio alone.
  \item Hormuz tanker capacity Granger-causes the DAP fertilizer
        price at a 1-month horizon ($p = 0.010$), with an IRF
        consistent in sign across all five layers.
  \item A small but consistent network contagion coefficient
        ($\rho \approx 0.11$) implies that direct shocks are
        amplified by $\sim\!13\,\%$ through bilateral co-movement.
\end{itemize}

\paragraph{Estimates that should be read with caution.}
The panel-regression coefficients $\beta_{\log p}$ and $\beta_{DWT}$
suffer near-collinearity with the time fixed effects (both regressors
are global, time-only series), so individual signs are not
identified --- only the country-varying interaction $\gamma_3$ is
robustly estimated. The DAP price stands in for ammonia, and
LNG-Japan for LPG butane; both are imperfect proxies and likely
understate spot-market transmission.

\paragraph{The wheat / ammonia ``positive'' result.}
The S2--S3 mean-percentage-change bars for wheat and ammonia are
slightly positive (Figure~\ref{fig:scenario-comparison}). This is an
artefact of the price channel: under our 2020--2024 estimation
window, log-prices and log-imports co-move positively (both rose
during the COVID--Ukraine shock years), and the SDM
$\beta_{\log p}$ inherits that correlation. We do not interpret it
as a structural elasticity; the structural channel for wheat is
near zero (Hormuz carries 0.1\,\% of wheat trade), so the residual
price effect dominates the average. The dollar impact on wheat
remains within the noise band of the estimation.

\paragraph{Other limitations.}
The simulation holds the 2020 trade matrix fixed and does not model
substitution. In reality, India and China would source LPG from the
US, Australia or West Africa at a price premium, so the
\emph{volume} losses we report are upper bounds and the realised
\emph{cost} losses (including substitution surcharges) may be higher
in dollar terms but distributed differently. Reserves data from
JODI could in principle inform a buffer-stock extension; we leave
that for future work.

% =====================================================================
\section{Conclusion}
% =====================================================================

This study quantifies the trade-network exposure of five commodity
layers to a Strait of Hormuz disruption, combining static
network-topology, community detection and targeted-attack analysis
with time-series and spatial econometrics, and integrates the two
into a transparent scenario simulation engine.

The empirical answer is unambiguous: \textbf{LPG is the systemic
vulnerability}. Wheat and most fertilizers can absorb a Hormuz
closure with modest re-routing; LPG cannot, because the Gulf
supplies nearly half of the world's seaborne butane and a quarter
of the world's propane. A 6-month closure scenario translates into
a \$10.2 billion annualised global LPG import loss, with India alone
accounting for three-quarters of the impact.

Two methodological lessons emerge. First, the trade-network
literature's standard topology metrics (density, clustering,
efficiency) miss the asymmetric exposure to a regional disruption
unless complemented by a node-removal experiment that respects
edge weights. Second, the spatial-econometric multiplier --- though
small in our marginal estimates --- compounds across countries and
should be a default amplifier in policy scenario engines.

% =====================================================================
\bibliographystyle{plainnat}
\bibliography{references}

\end{document}
